% 本文件由 LaTeX 排版 Agent 根据有效 translation.json 自动生成。
% author=Augustine of Hippo; work_id=1305; title=论未见之事的信德

\chapter{论未见之事的信德}

% structure: p0001 source=h1 target=omitted reason=page-title-already-rendered-by-parent

\EditorialNote{有些人认为这篇短论是伪作，但已知是圣奥斯定之作，因为他在……中提及} -->第231封书信 \Emph{ad Darium Comitem}。似乎写于399年之后，从关于偶像的第10节所述内容可知；因为在该年，霍诺留颁布了反对偶像的法律。—— \Emph{摘自本笃版}。\par

1. 有些人认为，基督宗教是应当加以嘲笑而非持守的，原因在于，其中所展示的不是可见之物，而是命令人相信未见之事的信德。因此，我们为要驳斥这些自以为明智而不愿相信他们所不能见之事的人，虽然我们不能将我们所信的神圣事物显于人眼，却仍能向人的心灵显示：即使是那些未见之事，也应当相信。首先，应当告诫他们（愚昧使他们如此屈从于肉体之眼，以致凡不能通过它们看见的，就以为不应相信），有多少事物他们不仅相信而且知道，却是这样的眼睛所不能见的。这些事物在我们自己的心灵中数不胜数（心灵的本性是无法看见的），更不必说其他了；正是那使我们相信的信德，或者使我们知道我们相信某事或不相信的念头，本身完全不属于那些眼睛的视野；对于我们的内在之眼，有什么如此赤裸、如此清晰、如此确定呢？那么，我们怎能不信我们用肉体之眼所不能见的事呢？然而，在不能运用肉体之眼的情况下，我们相信或不相信，这本身难道不是我们毫无疑惑地看见的吗？\par

2. 但他们说：那些在心灵中的事物，既然我们能凭心灵本身辨别它们，就无需通过肉体之眼来认识；但你们叫我们相信的那些事物，你们既不从外面指出来，使我们能通过肉体之眼认识它们；也不在我们的心灵之内，使我们能通过思考看见它们。他们这样说，仿佛如果有人已经能看见所信之物摆在面前，还会被要求去相信似的。因此，我们确实也应当相信某些未见之事的暂时事物，好使我们配得看见我们所信的永恒事物。但是，无论你是谁，若只相信你所看见的，看吧，眼前的身体你用肉体之眼看见，你自己眼前的意愿和思想，因为在你自己的心灵里，你凭心灵本身看见；请告诉我，你用什么眼睛看见你朋友对你的意愿？因为没有任何意愿能被肉体之眼看见。什么？你在自己的心灵里也能看见别人心里所发生的事吗？但如果你看不见，那么你如何回报你朋友的好意呢？如果你不相信你所不能看见的，你会回报吗？也许你会说，你通过朋友的作为看见他的意愿。因此你会看见行为，听见言语，但关于你朋友的意愿，那不能看见也不能听见的，你将相信。因为那意愿不是颜色或形状，可以投射在眼睛上；也不是声音或旋律，可以溜进耳朵里；它也不是你自己的，以致能通过你自己内心的运动被感知。因此，它既不被看见，也不被听见，也不被你自己内在地看见，就必须被相信，免得你的生命因没有友谊而荒芜，或者别人给予你的爱得不到你的回报。那么，你所说的\Quote{除了你从外面在身体上看见的，或者从内在在心里看见的，你不应当相信别的}这话又在哪里呢？看，你从你自己的心里，相信了一个不是你自己的心；你付出了你的信德，却没有将你的身体或心灵的目光投向他。你凭自己的身体辨认朋友的面容，你凭自己的心灵辨认自己的信德；但除非你里面有那样的信德，能相信在他身上你看不见的东西，否则你不会爱朋友的信德。尽管一个人也可能通过假装好意和隐藏恶意来欺骗；或者，如果他无意加害，却因期望从你那里得到好处而假装，因为他并没有爱。\par

3. 但你说，你所以相信你的朋友（你看不见他的心），是因为你在试炼中证实了他，也知道了他在你的危险中对你怀有怎样的心，他没有离弃你。那么，你觉得我们必须渴望自己的苦难，好使朋友对我们的爱得到证明吗？难道除非一个人因逆境而不幸，否则就不能拥有最可靠的朋友吗？如此，他若不是被自己的痛苦和恐惧所折磨，就不能享受对方经过考验的爱？在拥有真正朋友这件事上，怎么能渴望那种只有不幸才能验证的幸福，而不更应畏惧它呢？然而，确实，朋友在顺境中也同样可以拥有，但在逆境中更确切地得到证明。但可以肯定，为了证明他，你也不会把自己投入危险之中，除非你相信；因此，当你为了证明而投身时，你在证明之前就已经相信了。因为，如果我们不应相信未见之事，我们却相信朋友的心，而且那心还未确定地证明；并且，在我们通过自己的不幸证明他们为好之后，我们仍然相信他们对我们怀有好意，而非看见：只是信德如此之大，以至于我们并不不恰当地认为，我们以信德的某种眼睛看见了我们所相信的，然而我们应当相信，正因我们不能看见。\par

4. 如果这一信德从人间事务中被除去，有谁不会看到它们中间有多大的失序，以及多么可怕的混乱必将随之而来呢？因为，若我所看不见的，我本不应相信，那么谁会因相互的情感而被任何人爱呢？（因为爱本身是不可见的。）因此，整个友谊将会消亡，因为它只由相互的爱构成。因为，如果没有人相信任何爱的表露，它还能从任何人那里得到什么呢？再者，友谊消亡，心灵中既不会保存婚姻的纽带，也不会保存亲属和亲戚的纽带；因为在这些关系中，也肯定有一种情感上的友好结合。因此，配偶将无法相互爱慕，因为每方都不相信另一方的爱，因为爱本身不能被看见。他们也不会渴望生养儿子，因为他们不相信儿子会回报他们。如果这些孩子出生并长大，父母自己就更不会爱自己的儿女了，因为他们看不见儿女心中对自己的爱，这爱是不可见的；如果相信未见之事不是值得称赞的信德，而是应受指责的轻率。我现在为什么要谈论其他的关系，如兄弟、姐妹、女婿、岳父，以及所有因血缘或姻亲而联结的人呢？如果爱是不确定的，意志是可疑的——儿子对父母的爱，父母对儿子的爱——而应有的善意得不到回报；因为人们既不认为那是应当的，也不认为别人身上那不可见的东西是存在的。再者，如果这种谨慎不是才智的标志，而是可憎的——即我们不相信自己被爱，因为我们看不见爱我们之人的爱，也不回报那些我们认为不欠他们回报的人——那么，如果我们不相信所看不见的，人间事务就会陷入如此程度的混乱，以致完全彻底地颠覆，如果我们不相信任何人的意志（我们当然看不见它们）。我省略不谈，那些责备我们相信未见之事的人，他们在多少事情上相信传闻或历史，或者相信他们自己从未到过的地方；他们并不说：\Quote{我们不相信，因为我们没有看见。}因为如果他们这样说，他们就不得不承认，他们自己的父母对他们来说不是确实知道的：因为在这一点上，他们也相信了别人讲述的叙述，而这些人无法显示它，因为这是一件已经过去的事；他们自己对那个时代没有感觉，却毫无疑虑地赞同那个时代的其他人的说法：如果不这样做，就必然会对父母犯下不虔不孝的罪过，而我们似乎是在表明对那些我们看不见的事情相信的轻率。因此，既然如果我们不相信那些我们看不见的事物，人类社会本身就会因和谐消亡而无法存续，那么信德岂不更应当应用于神圣的事物上，尽管它们看不见？如果缺乏这种应用，被违反的就不是某些人的友谊，而是最首要的虔诚的纽带，从而招致最大的痛苦。\par

5. 但你会说：朋友对我的善意，我虽不能看见，却能通过许多证据推究出来；但你，你要我们相信那些看不见的事，却没有证据可以显示它们。同时，你承认因某些证据的显明，有些事即使看不见也应当相信，这并非小事：因为由此达成共识，并非所有看不见的事都不该相信；而那句\Quote{我们不应相信看不见的事}的话，便被驳倒、丢弃、推翻。但如果有人认为我们相信基督，却没有关于基督的任何证据，那他们大错特错了。因为有什么证据，比我们现在看到那些预言并应验的事更清楚呢？因此，你们这些认为没有理由相信那些关于基督的未见之事的人，请留意你们所看见的事。教会本身以慈母之口向你们说话：\Quote{我在普世受你们惊叹，结出果实并增长，昔日并非如你们现在所见。}但，\BibleQuote{因你的后裔，地上万民都要蒙受祝福。}当天主祝福亚巴郎时，祂应许了我；因为在基督的祝福中，我遍传于万民之中。基督是亚巴郎的后裔，历代相传的次序为此作证。简言之，亚巴郎生了依撒格，依撒格生了雅各伯，雅各伯生了十二个儿子，以色列民族即由此而出。因为雅各伯本人被称为以色列。在这十二个儿子中，他生了犹大，犹太人由此得名，童贞玛利亚出自犹大支派，她生了基督。看啊，在基督内，即亚巴郎的后裔中，万民都蒙受祝福，你们看见并惊叹：你们还害怕相信祂吗？你们更该害怕不相信祂才是。什么？你们怀疑或拒绝相信一位童贞女的产痛吗？你们更该相信，天主这样降生为人是合宜的。因为你们也应当接受此事已由先知预言：\BibleQuote{看，一位童贞女将怀孕生子，人要称祂的名字为\OriginalTerm{厄玛努尔}{Emmanuel}，意思是：天主与我们同在。}因此，你若愿意相信一位天主的诞生，就不会怀疑一位童贞女的生产；祂并未放弃对世界的治理，却以肉身来到人间；使祂的母亲获得生育之果，却不损其童贞。因为祂虽永是天主，却必须这样降生为人，借此诞生祂可以成为我们的一位天主。为此，先知再次论及祂说：\BibleQuote{天主啊，你的宝座永远常存；你王国的权杖是正义的权杖。你喜爱正义，憎恶不义；因此天主，你的天主，用喜乐之油傅了你，胜过你的同伴。}这傅油是属神的，天主用此傅油了天主，即圣父傅了圣子；由此从\Quote{\OriginalTerm{傅油}{Chrism}}，即傅油，我们知道祂是基督。我是教会，在同一个圣咏中，论及教会并对祂说，未来之事被预言为已完成：\BibleQuote{王后佩戴金饰，穿着五彩锦衣，站在你的右边。}即在智慧之奥秘中，\Quote{以各种语言装饰。}那里对我说：女子，请听，请看，请侧耳，忘记你的民族和你的父家；因为君王恋慕你的美艳，祂是你的主天主；提洛的女儿要前来献礼，民间最富有的也要向你求恩。君王女儿的一切荣耀都在内里，有金线流苏，穿着五彩锦衣。她后面有童女陪伴，被带到君王面前；她们要欢欣鼓舞地被引入君王的宫殿。你的祖先要由你的子孙替代，你要立他们做普世的王侯。他们要代代纪念你的名字。因此万民要永远称谢你，直到永远。\par

6. 倘若你们看不到这位王后，如今还因王室后嗣而富足；倘若她看不到那曾应许给她的事得以成就，那位曾被告知\Quote{女儿，请听，请看}的她；倘若她尚未弃绝世界古老的仪式，那位曾被告知\Quote{忘却你的人民和你父的家室}的她；倘若她不随处宣认基督为主，那位曾被告知\Quote{君王恋慕你的美丽，因为他是你的主天主}的她；倘若她不见万城之民向基督倾吐祈祷、献上礼物，关于基督曾对她说\Quote{提洛的女儿要前来向你献礼}；倘若富人的骄傲也未曾放下，他们不向教会求助，关于教会曾被告知\Quote{民间的富豪也要恳求你}；倘若他不承认那位君王的女儿，她被吩咐向他说\Quote{我们在天上的父}；并且在她的圣徒内，内在的人不日日更新，关于她曾被告知\Quote{那位君王女儿的一切光荣都在内里}——尽管她也以外在的传道者的名声之光，在言语的多样性中，照耀那些外在之人的眼目，如同\Quote{金线绣成的衣袍和彩色锦衣}；倘若如今，因他的美名传遍各处，没有贞女被带来献给基督以行奉献礼，关于她们曾说到，也对她们说道\Quote{贞女们要随她之后被带到君王面前，她的同伴们也要被带到你面前}；并且，为使她们不显得像被俘虏带进某座监狱，他说\Quote{她们要在欢乐和喜乐中被带领，她们要被带进君王的圣殿}；倘若她不生育儿子，好从他们中得到如同父亲的人，她可以随处立他们为统治者，关于她曾被告知\Quote{你的儿子要代替你的祖先，你要立他们为王侯，统治天下}——他们的母亲既在祈祷中优先于他们又服从于他们，将自己托付给他们：\Quote{他们必世世代代怀念你的名}；倘若因这些父亲们的宣讲——他们不停提及她的名——在她内没有聚集如此众多的群众，且在他们自己的语言中无休止地宣认恩宠的赞美，关于她被说\Quote{因此万民要永远歌颂你，直到永远}；倘若这些事情没有如此清晰显明，以致仇敌的眼目找不到转向何处而不被同样的清晰所击中，从而被迫承认那显然之事：那么你们或许有理由断言，没有证据向你们展示，你们可以通过看到这些证据而相信那些你们看不见的事。但如果那些你们看见的事早已被预言且如此清晰地应验；如果真理本身通过先后的效果向你们显明自己，噢，不信的残余，为使你们相信那看不见的事，要为你们所看见的事而羞愧。\par

7. 教会向你们说：\Quote{你们要听我}，你们要听我，就是你们所看见的，虽然你们不愿意看见。因为当时在\Place{犹太}地区的信友们，亲身在场并学习了\Person{基督}由童贞女奇妙的诞生、祂的受难、复活、升天，以及祂所有神圣的言行。这些事你们没有看见，因此你们拒绝相信。所以要注视这些事，定睛在这些事上，反思你们所看见的，这些事不是作为过去之事告诉你们，也不是作为未来之事预言给你们，而是作为现在之事显示给你们。什么？你们以为这是虚妄或微小的事，且认为这不是或只是一个小的神圣奇迹，即因被钉者的圣名，全人类都奔涌而来？你们没有看见那关于\Person{基督}人性诞生所预言并应验的事：\BibleQuote{看，童贞女将怀孕生子，}但你们看见了那向\Person{亚巴郎}所预言并应验的天主圣言：\BibleQuote{地上万民都要因你的后裔蒙受祝福。}你们没有看见那关于\Person{基督}奇事所预言的：\BibleQuote{你们来，观看上主的作为，祂在地上行了何等奇事，}但你们看见了那所预言的：\BibleQuote{上主对我说：你是我的儿子，我今日生了你；你向我请求，我必将万民赐你作产业，将地极赐你为基业。}你们没有看见那关于\Person{基督}受难所预言并应验的事：\BibleQuote{他们刺穿了我的手和我的脚，数算了我所有的骨头；但他们自己却注视并观看我；他们分了我的衣服，为我的长衣拈阄。}但你们看见了在同一篇圣咏中所预言、如今明显应验的事：\BibleQuote{大地四极都要记起，并归向上主，万邦的各族都将在祂面前朝拜；因为王国属于上主，祂将治理万邦。}你们没有看见那关于\Person{基督}复活所预言并应验的事，圣咏以祂的身份说话，首先关于出卖祂的人和迫害者：\BibleQuote{他们走出门外，一同谈论：我的所有仇敌私下议论我，他们图谋恶事加于我；}他们用不义之言攻击我。为表明他们杀害那将要复活者毫无所获，祂接着说：\BibleQuote{怎么？那睡觉的，岂不还要起来吗？}稍后，在同一预言中关于那出卖者自己，祂预言了福音中所记载的：\BibleQuote{那吃过我饭的人，竟用脚踢我，}即践踏我。祂立即接着说：\BibleQuote{但是，上主，求祢怜悯我，使我再起来，我将报复他们。}这事已经应验：\Person{基督}睡了又醒了，即复活了；祂通过同一预言在另一篇圣咏中说：\BibleQuote{我躺下安睡，我又醒了，因为上主扶持了我。}但你们没有看见这些，却看见了祂的教会，关于她同样有记载并已应验：\BibleQuote{上主我的天主，万民要从地极来到祢面前，并要说：我们的祖先实在敬拜了虚假的偶像，其中毫无益处。}这一点，无论你们愿意与否，你们确实看见了；尽管你们仍然相信那些偶像目前或曾经有某种益处，但确实无数的外邦民族，在离开、抛弃或打碎诸如此类的虚妄之后，你们听见他们说：\BibleQuote{我们的祖先实在敬拜了虚假的偶像，其中毫无益处；人岂能制造神祇？看，它们并不是神！}不要以为预言论及外邦人要来到天主的某个地方，因为经上说：\BibleQuote{万民要来到祢面前，从地极。}你们如果能明白，就应知道，外邦民族来到基督徒的天主——至高的真天主——面前，不是靠行走，而是靠相信。因为同样的事由另一位先知如此预言：\BibleQuote{上主必向他们显现，彻底消灭地上万邦的一切神祇；列邦岛屿的人各从本地朝拜祂。}而一位说：\BibleQuote{万民都要来到祢面前，}另一位却说：\BibleQuote{他们各从本地朝拜祂。}所以，他们来到祂面前，并未离开自己的地方，因为因着相信祂，他们在心中找到了祂。你们没有看见那关于\Person{基督}升天所预言并应验的事：\BibleQuote{天主，请祢高居于诸天之上，}但你们看见了紧接着的下文：\BibleQuote{愿祢的荣耀充满大地。}那些关于\Person{基督}已经完成和过去的事，你们全都未看见；但祂教会中这些现在之事，你们不能否认你们看见。这两件事我们都指出是预言过的；但两件事的应验我们却不能指给你们看，因为我们无法使过去的事重现眼前。\par

8. 但就如朋友的意愿（虽不可见）透过可见的记号而得到相信；同样，现今可见的教会，对于一切不可见之事——那些在其自身亦被预言的著作中所显示的事——乃是过去的指标与未来的宣告。因为过去的事（现已无法看见）和未来的事（当时尚不可见），在它们被预言的时候，没有一件是当时可见的。因此，既然它们已开始按预言成就，从已成就的事到将要成就的事，那些关于基督和教会的预言已按顺序延续；属于这系列的，有关于审判之日、死者复活、不敬虔者与魔鬼一同永罚，以及敬虔者与基督一同永赏的事，这些事同样已被预言，却尚未到来。那么，当我们拥有作为二者见证的中间事物（我们看见的），并在先知书中听到或读到关于起初之事、中间之事和末后之事都在它们发生前被预言了，我们为什么还不相信那看不见的起初与末后之事呢？除非，也许不信的人断定这些事是基督徒所写，为使他们已相信的事（若被认为是在发生前应许的）具有更大的权威分量。\par

9. 若他们怀疑这一点，就让他们仔细查考我们敌人犹太人的抄本。在那里，让他们读到我们所提及的、关于我们所信的基督和关于我们从信德艰辛开端以至王国永恒福乐中所辨明的教会而预言的那些事。但他们在阅读时，不要因拥有这些书卷的人因仇恨的黑暗而不明白而感到奇怪。因为他们不会明白，这早已被同一些先知预言；这应该像其余的一样应验，并且藉由天主隐秘而公义的审判，应得的惩罚将归于他们的行为。确实，祂——他们钉死并在木架上给祂苦胆和醋的那位——虽因那些祂即将从黑暗中领到光中的人，曾向圣父说：\BibleQuote{赦免他们，因为他们不知道他们所做的是什么；}然而因那些祂因更隐密的原因即将舍弃的人，早已由先知如此长久地预言：\BibleQuote{他们拿苦胆给我作食物，我渴时，他们拿醋给我喝；愿他们的筵席在他们面前成为网罗、报应和绊脚石；愿他们的眼睛昏花，不得看见，并使他们的腰时常弯曲。}这样，他们虽有我们事业最清晰的见证，却以昏花的眼睛四处行走，为的是藉着他们，那些他们自身被否定的见证得以证实。因此，这事得以成就，使他们不至于被完全涂抹，以致这个教派全然不存在；而是被分散在全地，为的是带着赐予我们的恩宠的预言，更确实地说服不信者，处处使我们受益。而我所断言的这件事，请你们接受它是如何被预言的：祂说：\BibleQuote{不可杀他们，恐怕他们忘记祢的法律，要用祢的能力分散他们。}因此他们没有被杀，因为他们没有忘记在他们中间所诵读和所听见的事。因为若他们完全忘记——尽管他们不明白——圣经，他们就会在犹太教仪式本身中被杀；因为当犹太人不再知道法律和先知时，他们就不能使我们受益。因此，他们没有被杀，而是被分散；为的是，虽然他们在信德上不能得救，却应在记忆中保留那些可以帮助我们的东西；在他们的书中，是我们的支持者；在他们的心里，是我们的仇敌；在他们的抄本中，是我们的见证人。\par

10. 虽然，即使关于基督和教会的见证并未先前存在，谁不应被这事实所感动而相信呢：神圣之光明突然照耀人类，当我们看到（假神已被抛弃，他们的偶像到处被打碎，他们的庙宇被推翻或改为他用，这么多虚妄的仪式从人类最根深蒂固的习惯中被连根拔起）所有人都呼求唯一真天主？而这事是由一人成就的——那人被嘲笑、被抓住、被捆绑、被鞭打、被掌掴、被辱骂、被钉十字架、被杀害。他的门徒（他拣选了普通人、无学问的人、渔夫和税吏，借他们来传扬他的教导）宣扬他的复活、他的升天，他们声称亲眼看见了这些，并且充满圣神，用他们从未学过的各种语言大声宣讲这福音。听他们的人中，一部分相信了，一部分不信，激烈地反对传道者。这样，他们为真理至死忠心，不以恶报恶，而是忍受；不是通过杀人而是通过死去来胜过；这样，世界被改变归向这宗教，这样，凡人的心转向这福音——男人女人、小的大的、有学问的无学问的、智慧的和愚拙的、有能力的软弱的、尊贵的卑贱的、高的低的——并且教会传遍万邦，如此增长，以至于甚至针对天主教信仰本身，没有任何异端邪说、任何错误，被发现如此反对基督真理，以至于它不试图以基督之名夸耀；这种错误本身不会在地上兴起，除非这反对本身实施了一种有益的管教。那位被钉十字架的怎能如此大能，假如他不是成为人的天主，即使他没有借先知预言这些事呢？但现在这如此伟大的虔敬奥秘已经有了它的先知和先驱先行宣告，借他们的神圣声音预先宣布了；并且它按预先宣告的方式来到，谁还如此疯狂，竟断言宗徒们关于基督说了谎？他们宣讲基督是按先知预先宣告他要来的方式而来的，那些先知关于宗徒们自己也没有对将要来的真事缄默。因为关于这些，他们曾说过：\BibleQuote{他们的言语没有可听的，他们的声音传遍全地，他们的话语传到地极。}而这我们至少看到在世界中实现了，尽管我们尚未在肉身中看见基督。因此，除非被惊人的疯狂弄瞎，或被惊人的顽固所硬化，谁不愿相信圣经呢？圣经预言了全世界的信仰。\par

11. 但是，亲爱的朋友们，你们拥有这信德，或现在刚开始拥有它，愿它在你们内被滋养而增长。因为正如暂时的事物已经来到，早在许久之前被预言，同样永恒的事物也将来到，这些是所应许的。不要让虚妄的外邦人、虚假的犹太人、狡猾的异端者，或者甚至在天主教会内部的恶基督徒欺骗你们——后者是更有害的敌人，因他们越是在我们内。因为，唯恐软弱者在这事上受困扰，神圣的预言并未缄默，在\BibleBook{}{雅歌}中新郎对新娘说，即主基督对教会说：\BibleQuote{我的佳偶在女子中，如同百合花在荆棘中。}他没有说\Quote{在外人中间}，而是说：\BibleQuote{在女子中。凡有耳的，听吧！}并且，正如圣福音所说，撒在海里的网聚集各种鱼，正被拖到岸边，即到世界的末了，让他自己在心里（而非在肉体上）与恶鱼分离；通过改变恶习，而非打破圣网；以免那些现在看似被选的人与弃者混杂，当他们在岸边开始被分离时，发现不是永生，而是永罚。\par

\SourceRecord{Translated by C.L. Cornish.；Nicene and Post-Nicene Fathers, First Series；Vol. 3.；Edited by Philip Schaff.；Buffalo, NY: Christian Literature Publishing Co.,；1887.；<http://www.newadvent.org/fathers/1305.htm>.}
